\documentclass{article}
\usepackage[utf8]{inputenc}
\usepackage{graphicx}
\usepackage{amsmath}
\usepackage{amssymb}
\usepackage{geometry}
\usepackage{hyperref}
\usepackage{xurl}
\usepackage{cite}
\usepackage{indentfirst}
\usepackage{yhmath}
\usepackage{float}
\usepackage{listings}
\usepackage[dvipsnames]{xcolor}
\usepackage{booktabs} 
\usepackage{array}
\usepackage{titlesec}
\usepackage{multirow, multicol}
\geometry{letterpaper, margin=1in}
\hypersetup{hidelinks}
\def\UrlBreaks{%
    \do\/%
    \do\a\do\b\do\c\do\d\do\e\do\f\do\g\do\h\do\i\do\j\do\k\do\l%
    \do\m\do\n\do\o\do\p\do\q\do\r\do\s\do\t\do\u\do\v\do\w\do\x\do\y\do\z%
    \do\A\do\B\do\C\do\D\do\E\do\F\do\G\do\H\do\I\do\J\do\K\do\L%
    \do\M\do\N\do\O\do\P\do\Q\do\R\do\S\do\T\do\U\do\V\do\W\do\X\do\Y\do\Z%
    \do0\do1\do2\do3\do4\do5\do6\do7\do8\do9\do=\do/\do.\do:%
    \do\*\do\-\do\~\do\'\do\"\do\-}
\urlstyle{same}

\title{\textbf{ECE3700J Introduction to Computer Organization \\
Lab 4 Report --- Resolving Hazards}}

\author{\large Group 12 \\ % Update with your names maybe
% \large Student 1 Name (ID) \\
% \large Student 2 Name (ID) \\
% \large Student 3 Name (ID)
}

\date{\today} % Or \date{November 15, 2025}

\lstset{
    language=Verilog,
    backgroundcolor=\color{gray!20},
    basicstyle=\ttfamily\small,
    breakatwhitespace=false,
    breaklines=true,
    captionpos=b,
    keepspaces=true,
    numbers=left,
    numbersep=5pt,
    showspaces=false,
    showstringspaces=false,
    showtabs=false,
    tabsize=2,
    frame=single,
    framerule=0pt,
    framesep=3mm,
    xleftmargin=3mm,
    xrightmargin=3mm
}

\begin{document}

\maketitle

\newpage
\tableofcontents

\newpage
\section{Introduction}

This lab continues from the single-cycle processor of Lab 2 and the pipelined core of Lab 3 to implement a 32-bit, 5-stage pipelined RISC-V processor with full hazard support. The design is implemented in synthesizable Verilog HDL.

The processor supports an expanded subset of the RV32I instruction set, including:
\begin{itemize}
    \item \textbf{Arithmetic/Logical (R-type):} \texttt{add}, \texttt{sub}, \texttt{and}, \texttt{or}, \texttt{sll}, \texttt{srl}, \texttt{sra}
    \item \textbf{Arithmetic/Logical (I-type):} \texttt{addi}, \texttt{andi}, \texttt{slli}, \texttt{srli}, \texttt{srai}
    \item \textbf{Memory Access:} \texttt{lw}, \texttt{sw}, \texttt{lb}, \texttt{lbu}, \texttt{sb}
    \item \textbf{Conditional Branch:} \texttt{beq}, \texttt{bne}, \texttt{bge}, \texttt{blt}
    \item \textbf{Unconditional Jump:} \texttt{jal}, \texttt{jalr}
\end{itemize}

The primary goal of Lab 4 is to extend the Lab 3 pipelined core with hardware support for data and control hazards. We add a dedicated \texttt{forwarding\_unit} in the EX stage to resolve ALU--ALU and ALU--store hazards, a \texttt{hazard\_detection} unit in ID to handle load--use hazards via one-cycle stalls, and control-logic to perform early branch/jump resolution in ID with IF/ID flushing. The correctness of these mechanisms is verified using the official \texttt{Lab4\_testcase.s} program and waveform inspection in Vivado.

\section{Processor Architecture}

The processor implements a classic 5-stage RISC-V pipeline: Instruction Fetch (IF), Instruction Decode (ID), Execute (EX), Memory (MEM), and Write Back (WB). Four pipeline registers are used to pass data and control signals between stages: \texttt{if\_id\_reg}, \texttt{id\_ex\_reg}, \texttt{ex\_mem\_reg}, and \texttt{mem\_wb\_reg}.

\subsection{IF Stage }

The IF stage contains the \texttt{pc\_reg}, which updates the program counter only when \texttt{PC\_write} is asserted. This allows the hazard detection unit to freeze the IF stage during load–use hazards. Each cycle, the instruction memory is accessed using the current \texttt{pc}, and the fetched instruction is stored into \texttt{if\_id\_reg}, unless \texttt{IF\_ID\_write} is deasserted (stall) or \texttt{IF\_Flush} is asserted (control-hazard flush).

Unlike a conventional design where the branch target is computed in EX, our implementation moves all branch/jump PC-selection logic into the ID stage. A 3-way multiplexer in ID selects:
\begin{enumerate}
    \item \texttt{PC + 4} (default sequential execution),
    \item \texttt{PC + imm} (for taken branches or \texttt{jal}),
    \item \texttt{rs1 + imm} with LSB cleared (for \texttt{jalr}).
\end{enumerate}
The selected value becomes \texttt{q\_pc\_new}, the next PC fed into \texttt{pc\_reg}.
This early-PC-resolution strategy minimizes control hazards to a single-cycle flush.


\subsection{ID Stage}

In ID, the \texttt{control} unit decodes the instruction, and the \texttt{reg\_file} simultaneously reads \texttt{rs1} and \texttt{rs2}. The immediate is generated by \texttt{imm\_gen}. All these signals are forwarded to the ID/EX register unless a stall is inserted.

\textbf{Hazard Detection (Load–Use Stall).}
The \texttt{hazard\_detection} module checks whether:
\[
\texttt{id\_ex\_mem\_read = 1} \quad \text{and} \quad \texttt{id\_ex\_rd} = \texttt{rs1} \text{ or } \texttt{rs2}.
\]
If so, it asserts:
\begin{itemize}
    \item \texttt{PC\_write = 0}: freeze IF,
    \item \texttt{IF\_ID\_write = 0}: freeze ID,
    \item \texttt{clear\_control = 1}: insert a bubble into ID/EX.
\end{itemize}

\textbf{Control Hazard Resolution.}
Branch comparison and jump target computation are performed in ID using a dedicated \texttt{compare} module and a \texttt{check\_sel} unit. Forwarding from EX/MEM to ID ensures that the branch uses the most recent operand values. When a branch or jump is taken:
\begin{itemize}
    \item \texttt{IF\_Flush = 1} clears the IF/ID register,
    \item \texttt{q\_pc\_new} is updated to the branch/jump target.
\end{itemize}
This ensures correct control-flow redirection with only a single-cycle penalty.


\subsection{EX Stage}

In EX, the pipeline resolves data hazards via a structured forwarding network. The \texttt{forwarding\_unit} compares:
\[
\texttt{ex\_mem\_rd}, \quad \texttt{wb\_rd}, \quad \texttt{id\_ex\_rs1\_addr}, \quad \texttt{id\_ex\_rs2\_addr}
\]
and generates \texttt{forwardA} and \texttt{forwardB}. Two 3-way multiplexers select the correct operands for the ALU from:
\begin{enumerate}
    \item ID/EX register file outputs,
    \item MEM/WB write-back value,
    \item EX/MEM ALU result.
\end{enumerate}

This eliminates ALU–ALU and MEM–ALU hazards without requiring pipeline stalls. The ALU operation is determined by \texttt{alu\_control}. The result, together with forwarded store data and relevant control fields, is passed into the EX/MEM register.


\subsection{MEM Stage}

The MEM stage performs load and store operations through \texttt{data\_mem}. A small MEM-hazard fix is included: if a store immediately follows a load that writes to the same register used as its store-data input, the design forwards the just-loaded value from MEM/WB into the store-data path.

The outputs of this stage—either loaded data or the ALU result—are stored in the MEM/WB register for final write-back.


\subsection{WB Stage}

The WB stage selects the appropriate value to write into the register file:
\begin{itemize}
    \item ALU result (arithmetic instructions),
    \item memory data (loads),
    \item \texttt{PC+4} (for \texttt{jal} and \texttt{jalr} return address).
\end{itemize}
Write-back is enabled only when \texttt{RegWrite} is asserted in MEM/WB. This concludes the pipeline execution for one instruction while the next ones continue through earlier stages.


\section{Hazard Handling and Simulation Results}

In this section, we present how the pipeline resolves data and control hazards, using waveform snapshots captured from running the official \texttt{Lab4\_testcase.s} program. Each simulation figure corresponds to a specific instruction sequence where a hazard is expected to occur. The waveforms highlight the key control signals, forwarding paths, stalls, and pipeline flushes, demonstrating that our implementations of data hazard forwarding, load–use stall insertion, and control hazard flush behave correctly during execution.

\subsection{EX/MEM Data Hazards (Forwarding)}

Figure~\ref{fig:ex_mem_hazard} shows the full simulation of \texttt{Lab4\_testcase.s}.
In the region where \texttt{addi t1, x0, 0x399} is followed by \texttt{sw t1, 4(x0)},
signal \texttt{fwdB} is driven to \texttt{2'b10} and \texttt{mem\_wdata} becomes
\texttt{0x00000399}, indicating that the store data is forwarded directly from
the EX/MEM ALU result (no stall and the correct value is written to memory).
Later, for the chain \texttt{add t2, t0, t3} $\rightarrow$ \texttt{and t1, t2, t3}
$\rightarrow$ \texttt{andi t1, t2, 0} $\rightarrow$ \texttt{sub t0, t1, x0}, the
waveform shows \texttt{fwdA}/\texttt{fwdB} selecting EX/MEM on the first dependent
instruction and MEM/WB on the second, while the PC continues to advance;
the final register values match the expected results, confirming that both
EX and MEM data hazards are correctly resolved by the forwarding unit.

\begin{figure}[H]
    \centering
    \includegraphics[width=\linewidth]{Figures/sim_ex_mem_full.jpeg}
    \caption{Overall simulation including EX/MEM data hazards
    (highlighted windows show \texttt{addi}~$\rightarrow$~\texttt{sw} and
    the \texttt{add}/\texttt{and}/\texttt{andi}/\texttt{sub} chain).}
    \label{fig:ex_mem_hazard}
\end{figure}

\subsection{\texttt{load-use} Hazard}

\textbf{Explanation:}
This waveform shows the simulation of a \texttt{lw} instruction (PC = 14) followed by a \texttt{bne} instruction (PC = 18), which together create a classic \texttt{load–use} hazard. As the \texttt{bne} instruction enters the ID stage at \texttt{15.000 ns}, the hazard detection unit asserts \texttt{clear\_control} and deasserts both \texttt{PC\_write} and \texttt{IF\_ID\_write}, inserting a one-cycle stall into the pipeline. 

After the \texttt{lw} instruction reaches the MEM stage, \texttt{br\_op2\_id} correctly receives the loaded value \texttt{0xffffff99} through the MEM-to-ID forwarding path, allowing the branch comparison to proceed with the updated operand. The waveform confirms that both the hazard detection logic and the stall insertion mechanism function as intended.



\begin{figure}[H]
    \centering
    \includegraphics[width=0.9\linewidth]{Figures/load-use.png} 
    \caption{Simulation of \texttt{lw t3 0(x0)}  \texttt{bne t0 t3 wrong -branch} (load-use hazard)}
    \label{fig:jal_sim}
\end{figure}

\subsection{\texttt{load-store} Hazard}

\textbf{Explanation:}
This waveform shows the simulation of the instruction sequence
\texttt{sw x5,0(x1)} (PC = 8), \texttt{lw x6,0(x1)} (PC = 12),
\texttt{sw x6,4(x1)} (PC = 16), and \texttt{lw x7,4(x1)} (PC = 20) from the
test program. Because \texttt{sw x6,4(x1)} immediately follows
\texttt{lw x6,0(x1)}, the value to be written by the store depends on the
result of the preceding load, creating a classic \texttt{load–store} hazard
on register \texttt{x6}.

The signals \texttt{dm\_word0[31:0]} and \texttt{dm\_word4[31:0]} monitor
the data memory words at addresses \texttt{0} and \texttt{4}, respectively.
Initially, both locations read \texttt{0x00000000}. After the
\texttt{sw x5,0(x1)} instruction reaches the MEM stage, 
\texttt{dm\_word0} is updated to \texttt{0x000001a2}, confirming that the
first store writes the test pattern correctly. A few cycles later, once
\texttt{lw x6,0(x1)} has loaded the value \texttt{0x000001a2} into
\texttt{x6} and \texttt{sw x6,4(x1)} executes, \texttt{dm\_word4} changes
from \texttt{0x00000000} to the same value \texttt{0x000001a2}. The fact
that location \texttt{4} receives the newly loaded value rather than the
old value \texttt{0x00000000} demonstrates that the processor correctly
resolves the \texttt{load–store} hazard (via forwarding and/or stalling),
and therefore the hazard-handling mechanism operates as intended.

\begin{figure}[H]
    \centering
    \includegraphics[width=0.9\linewidth]{Figures/load-store.png} 
    \caption{Simulation of load-store hazard}
    \label{fig:lw-sw-sim}
\end{figure}

\subsection{\texttt{control} Hazard}

\textbf{Explanation:}
This waveform shows the execution of an \texttt{andi} instruction (PC = 0x24), a \texttt{sub} instruction (PC = 0x28), and the subsequent \texttt{bge} instruction (PC = 0x2C). The \texttt{bge} depends on the results produced by the previous two ALU instructions, and it also introduces a control hazard since the target PC must be resolved in the ID stage.

When the \texttt{bge} instruction enters the ID stage at \texttt{29 ns}, the pipeline first handles the EX hazard between \texttt{andi} and \texttt{sub} by asserting \texttt{clear\_control} and deasserting both \texttt{PC\_write} and \texttt{IF\_ID\_write}, inserting a single-cycle bubble into the ID/EX register to avoid using stale operands.

Once the data hazards are resolved, the ID-stage branch comparison determines that the \texttt{bge} condition is satisfied. The control unit then asserts \texttt{IF\_Flush}, clearing the IF/ID register and redirecting the PC to the branch target at \texttt{0x34}. The final value of \texttt{q\_pc\_new} confirms that the pipeline correctly performs early branch resolution and uses a single-cycle flush to handle the control hazard.

\begin{figure}[H]
    \centering
    \includegraphics[width=0.9\linewidth]{Figures/control_hazard.png}
    \caption{Simulation waveform for the \texttt{andi} → \texttt{sub} → \texttt{bge} sequence, illustrating one-cycle stall insertion and IF/ID flush during branch resolution.}
    \label{fig:control_hazard}
\end{figure}



\subsection{\texttt{jal}/\texttt{jalr} Instructions}

\textbf{Explanation:}
This waveform illustrates the execution of a \texttt{jal x1, jump\_test} instruction at PC = 0x34. Since \texttt{jal} modifies the program counter, it introduces a control hazard similar to a taken branch. When the \texttt{jal} instruction reaches the ID stage at \texttt{35 ns}, the ID-stage jump logic immediately computes the correct target PC (0x40). 

To prevent the wrong-path instruction at PC = 0x38 from entering the pipeline, the control unit asserts \texttt{IF\_Flush}, clearing the IF/ID register and discarding the sequentially fetched instruction. Meanwhile, \texttt{PC} is redirected to the jump target, and \texttt{PC+4} is forwarded through the pipeline to be written into \texttt{x1} as the return address.

The waveform confirms that the pipeline correctly performs early jump resolution, applies a single-cycle flush, and updates the PC to \texttt{0x40}, demonstrating the correct implementation of \texttt{jal}/\texttt{jalr} handling.


\begin{figure}[H]
    \centering
    \includegraphics[width=0.9\linewidth]{Figures/jal.png}
    \caption{Simulation waveform for \texttt{jal x1, jump\_test} (PC = 0x34), showing IF/ID flush and redirection of the PC to the jump target at \texttt{0x40}.}
    \label{fig:jal_sim}
\end{figure}

\newpage
\section{RTL Schematic}

The Verilog design was synthesized using Xilinx Vivado to verify its synthesizability and examine the hardware structure. The resulting RTL schematic (Figure~\ref{fig:rtl}) shows the high-level blocks corresponding to the pipeline stages and major components like the \texttt{reg\_file}, \texttt{alu}, and \texttt{data\_mem}, connected by the pipeline registers.

\begin{figure}[H]
    \centering
    \includegraphics[width=\linewidth]{Figures/schematic.png} 
    \caption{Vivado-generated RTL schematic of the 5-Stage Pipelined Processor.}
    \label{fig:rtl}
\end{figure}

\section{Source Code Listing}
\label{sec:source}

The complete Verilog source code (including \texttt{top} design and all instantiated submodules) is provided in the submission \texttt{lab4.zip} on Canvas.

\section{Conclusion}

A 32-bit, 5-stage pipelined RISC-V processor was successfully designed, implemented, and verified in Verilog. The processor supports a broader range of instructions than in the previous lab, including arithmetic, memory, branch, and jump operations.

Key outcomes:
\begin{itemize}
    \item A 5-stage pipeline (IF, ID, EX, MEM, WB) was implemented, with pipeline registers passing data and control signals between stages.
    \item The datapath and control logic were designed to handle R, I, S, B, and J-type instruction formats.
    \item Simulation of key instructions (\texttt{add}, \texttt{addi}, \texttt{lw}, \texttt{sw}, \texttt{beq}, \texttt{jal}) confirmed correct data flow, control signal generation, and pipeline operation for each instruction type.
    \item The design was successfully synthesized in Vivado, confirming the RTL is synthesizable and producing a schematic that matches the intended pipelined architecture.
    \item Data and control hazards were resolved using a forwarding unit, a load--use hazard detection/stall mechanism, a load--store fix, and early branch/jump resolution with IF/ID flush.
\end{itemize}

This lab demonstrates the fundamental principles of pipelined processor design, showing how instruction execution can be overlapped to improve performance.

\newpage
\begin{thebibliography}{99} 
\bibitem{Lab3} ECE3700J Teaching Team. \texttt{T6 Pipelined Processor.pdf}, \texttt{Lab3.pdf}, \texttt{Lab3\_testcase.s}, and \texttt{Lab3\_testcase.txt}. Retrieved on 27 Oct 2025 from \hyperlink{site}{oc.sjtu.edu.cn}.
\bibitem{Lab4} ECE3700J Teaching Team. \texttt{T7 Data Hazards.pdf}, \texttt{T8 Control Hazards.pdf}, \texttt{Lab4.pdf}, and \texttt{Lab4\_testcase.s}. Retrieved on 15 Nov 2025 from \hyperlink{site}{oc.sjtu.edu.cn}.
\end{thebibliography}

\end{document}